%%%% CAPÍTULO 1 - INTRODUÇÃO
%%
%% Deve apresentar uma visão global da pesquisa, incluindo: breve histórico, importância e justificativa da escolha do tema,
%% delimitações do assunto, formulação de hipóteses e objetivos da pesquisa e estrutura do trabalho.

%% Título e rótulo de capítulo (rótulos não devem conter caracteres especiais, acentuados ou cedilha)
\chapter{Introdução}\label{cap:introducao}

Um texto curto apresentando o capítulo.

\caixa{Atenção}{Para utilizar esse template é obrigatória a leitura do conteúdo do arquivo \texttt{readme.md}, que está neste projeto!}

\section{Considerações iniciais}\label{sec:consideracoesIniciais}

As considerações iniciais compõem um texto curto e geral apresentando uma visão geral e sucinta do assunto principal relacionado ao trabalho e a inserção do objeto de pesquisa nesse assunto \cite{Moore:2000:CMC:333067.333074}.

Em relação ao assunto, o apresentado nesta seção pode estar relacionado a trabalhos de outros autores ou ao assunto que fornece a fundamentação (motivação) para o trabalho a ser desenvolvido. Se o assunto está relacionado a trabalhos de outros autores, a contribuição do trabalho é definida em relação ao que já foi pesquisado nesse assunto. Se o assunto será utilizado para embasamento do que será proposto, explicitar como o trabalho se insere nesse assunto. A contribuição pode, ainda, estar relacionada a uma necessidade de mercado ou a uma oportunidade decorrente de algum problema real para o qual se pretender propor uma solução. Nesse caso, o assunto fornece um contexto teórico de suporte para o problema e/ou a solução.

O importante nesta seção é deixar claro do que se trata o trabalho (assunto ou tema), identificar o objeto de pesquisa, como será encaminhada a solução (procedimento metodológico, tecnologias, ferramentas utilizadas) e o que se pretende ao final do trabalho, sem explicitar a solução e os resultados.

\caixa{Atenção}{As seções a seguir são sugestões, converse com o seu orientador para ver quais seções devem ter em seu trabalho!}

\section{Objetivos}\label{sec:objetivos}

Um texto curto\footnote{Teste de nota de rodapé 1.} apresentando a seção.

\subsection{Objetivo geral}\label{subsec:objetivoGeral}

Implementar um algoritmo de aproximação para o cálculo do coeficiente de clustering local em grafos que representam redes complexas, realizando estudos de comparações com algoritmos exatos utilizados para a computação de tal medida.

\subsection{Objetivos específicos}\label{subsec:objetivosEspecificos}

\begin{itemize}
\item Implementar algoritmo de aproximação para o cálculo do coeficiente de clustering local proposto por Lima et al. \cite{Lima2022} utilizando linguagem de programação Python.

\item Realizar revisão bibliográfica sobre a área de teoria de grafos e complexidade de amostra, e disponibilizar material em língua portuguesa para disseminação desse conhecimento científico.
\end{itemize}

\section{Justificativa}\label{sec:justificativa}

O coeficiente de clustering local foi inicialmente proposto com a intenção de determinar se um grafo qualquer possuía a propriedade de ser um grafo de mundo pequeno \cite{WattsSmallWorld1998}. Esses grafos possuem algumas características que os tornam muito importantes quando estudamos conceitos de redes sociais e fenômenos de contágio de doenças. Eles possuem alta coesão local [18], ou seja, os nós tendem a formar grupos ou clusters densamente conectados, e possuem baixa distância média [19], logo, a distância média entre dois nós quaisquer no grafo é relativamente pequena, mesmo que o número total de nós no grafo seja muito grande. Isso significa que a maioria dos nós pode ser alcançado a partir de qualquer outro nó através de um pequeno número de etapas.

Essas características são de extrema importância quando estudamos redes sociais ou redes de transporte, por exemplo. Isso porque existe uma eficiência de comunicação, pois a baixa distância média facilita a rápida disseminação de informações. Assim como também existe robustez e resiliência nessas estruturas, significando que falhas em parte da rede não necessariamente desconectam toda a rede, pois há caminhos alternativos dentro dos clusters.
Grafos de mundo pequeno são úteis para modelar fenômenos naturais e sociais como a propagação de doenças, a disseminação de ideias e tendências entre vários outros fenômenos [20]. Logo, existir um algoritmo que calcule essa métrica de forma muito aproximada dos exatos, porém em tempo muito menor faz-se muito importante.

Mais do que a própria importância do coeficiente de clustering local, uma das motivações para a implementação do algoritmo proposto por Lima et al. (2022) – além do próprio fato de ser um algoritmo novo sem nenhuma implementação realizada antes – é o diferencial pela utilização da teoria de complexidade de amostras em um algoritmo de amostragem. Uma das vantagens de se utilizar a métrica da VC-Dimension de teoria de complexidade de amostra é que é possível obter um limite superior para o tamanho da amostra que é mais restrito do que os fornecidos pelas técnicas padrão de amostragem de Hoeffding e limitante da união [21]. Além disso, esses limitantes clássicos da literatura para algoritmos de amostragem, como limitante da união, Chernoff e Chebyshev, limitam-se ao tamanho da entrada e não levam em conta a estrutura combinatória de cada entrada, como a VC-Dimension faz no contexto de complexidade de amostras [22].

Outro grande fator motivador deste trabalho é a falta de material científico em português dentro da área de teoria de grafos e teoria de complexidade de amostra, tornando o desenvolvimento do conhecimento

\section{Estrutura do trabalho}\label{sec:estruturaTrabalho}

A estrutura do trabalho contém uma relação dos capítulos e uma descrição sucinta do que cada um deles contém. Esta seção fornece uma visão geral do trabalho no sentido da sua estrutura em capítulos\footnote{Teste de nota de rodapé 2.}.

\caixa{Atenção}{O OverLeaf está demorando muito para compilar o modelo com o Capítulo de Exemplos, que explica como usar o LaTeX. Assim, esse capítulo foi removido (está comentado para não compilar), mas há um arquivo chamado \texttt{exemploPDF.pdf}, na raiz do projeto, que contém esse capítulo de exemplos!}
